\documentclass[10pt,a4paper]{amsart}
\usepackage[margin=19mm]{geometry}
\usepackage{amsmath,amssymb,mathtools}
\usepackage[scr=rsfs,cal=euler]{mathalfa}
\usepackage{xfrac}
\usepackage[unicode,hidelinks,bookmarks=false]{hyperref}
\newcommand{\ie}{i.e.\ }
\newcommand{\cf}{cf.\ }
\newcommand{\resp}{respectively\ }
\newcommand{\Subsection}{\subsection}
\providecommand{\nobreakdash}{\nobreak\hbox{-}}
\setlength{\emergencystretch}{2em}
\multlinegap0pt
\usepackage{enumerate}
\usepackage{amssymb}
\usepackage{float}
\newtheorem{lemma}{Lemma}
\title{Existence of Ground State and Excited Spinning $Q$-Vortex Solitons on Finite Domains}
\author{Caroline Brumelot, Luciano Medina}
\date{Source: arXiv:2602.07676v1 (CC BY 4.0)}
\begin{document}
\maketitle

\noindent\emph{Source and licence.} Existence of Ground State and Excited Spinning $Q$-Vortex Solitons on Finite Domains. Version arXiv:2602.07676v1 (CC BY 4.0), distributed by arXiv under the Creative Commons Attribution 4.0 International licence, http://creativecommons.org/licenses/by/4.0/, as recorded in the arXiv licence metadata for this exact version and shown on the abstract page https://arxiv.org/abs/2602.07676v1. Full licence notice: \texttt{licenses/arxiv-2602.07676-CC-BY-4.0.txt}.\par\smallskip

\noindent\textit{Editorial scope.} Counted unit: the article's lemma neccCond3 (source0.tex lines 370--380). The article's complete original proof of it is delivered with it (source0.tex lines 382--426, one \texttt{proof} environment of 45 lines). No mathematics is written, repaired or re-derived: the statement and the proof are the article's own lines, in the article's own notation, and no retained line is substituted or re-flowed.  No further source-local context is needed: every cross-reference of every retained line resolves inside the two delivered spans.  Nothing was omitted from any retained span, so the package carries no omission record.  No retained line contains a citation, so the wrapper carries no bibliography and no bibliography entry.  No retained line contains a figure environment, an image-inclusion command or drawing code, so no image is delivered and the package declares no asset dependency.  Editorial formatting only, in addition: an AMS \texttt{amsart} preamble that restates the article's own declarations and macros used by the retained lines, namely the environments \texttt{lemma} on separate counters, together with the standard packages those lines need.  The retained environments print in this document as Lemma 1 in order of appearance, whereas the article prints the same items with the numbering of its own full document.  The complete native source of the article as distributed in the arXiv e-print for this exact version is bundled unchanged as \texttt{sources/194198/source0.tex}.
	\begin{lemma}
		Let $a,b,\lambda > 0$ such that $b > a^2/4$. For any $\omega$ satisfying
		\begin{align}
			2\lambda\left(b-\dfrac{1}{4}a^2\right)<\omega^2, \label{neccCond3}
		\end{align}
		there exists a $\phi_0\in H$ such that $I_{\omega}(\phi_0)<0$, provided $P > P^*$. The bound $P^*$ is given explicitly by
		\begin{align}
			P^* = \frac{\mathcal{B} + \sqrt{\mathcal{B}^2 + 4 \mathcal{A} \mathcal{C}}}{2\mathcal{A}},
		\end{align}
		where $\mathcal{A}, \mathcal{B}, \mathcal{C}$ are positive constants defined in the proof below.
	\end{lemma}

	\begin{proof}
		We construct a test function $\phi_0$ to verify the condition $I_{\omega}(\phi_0) < 0$. Let $t > 0$ be a parameter to be fixed later and define the trapezoidal function:
		\begin{align}
			\phi_0(\rho)=\left\{
			\begin{array}{cc}
				t\rho, &0\leq \rho < 1,\\
				t, &1\leq \rho \leq P-1,\\
				t(P-\rho), &P-1 < \rho\leq P.
			\end{array}\right.
		\end{align}
		First, we verify that $\phi_0 \in H$. The function $\phi_0$ is continuous on $[0, P]$ as the piecewise definitions agree at the transition points $\rho=1$ and $\rho=P-1$. Furthermore, $\phi_0(0)=0$ and $\phi_0(P)=0$, satisfying the boundary conditions. The weak derivative $\phi_{0,\rho}$ exists and is piecewise constant (taking values $t, 0, -t$), hence $\phi_{0,\rho} \in L^2(0, P)$. Thus, $\phi_0 \in H$.
		
		We now compute the action $I_{\omega}(\phi_0)$. The potential terms are dominated by the plateau region $[1, P-1]$. Integrating exactly, we find:
		\begin{align}
			\int_0^P \rho \phi_0^2 d\rho &= t^2\left(\frac{P^2}{2} - P + \frac{1}{3}\right), \\
			\int_0^P \rho \phi_0^4 d\rho &= t^4\left(\frac{P^2}{2} - P + \frac{1}{5}\right), \\
			\int_0^P \rho \phi_0^6 d\rho &= t^6\left(\frac{P^2}{2} - P + \frac{1}{7}\right).
		\end{align}
		The gradient term evaluates to $\int_0^P \rho \phi_{0,\rho}^2 d\rho = t^2 P$. For the centrifugal term, we use the upper bound valid for $P \ge 2$:
		\begin{align}
			\int_0^P \frac{\phi_0^2}{\rho} d\rho \le \int_0^1 \rho t^2 d\rho + \int_1^P \frac{t^2}{\rho} d\rho = \frac{t^2}{2} + t^2 \ln P.
		\end{align}
		Substituting these into $I_{\omega}(\phi_0)$ and grouping by powers of $P$, we obtain:
		\begin{align}
			I_{\omega}(\phi_0) \le -\mathcal{A} P^2 + \mathcal{B} P + \mathcal{C} \ln P,
		\end{align}
		where the coefficient of the quadratic term is:
		\begin{align}
			\mathcal{A} = -\frac{\lambda}{2} \left[ t^6 - a t^4 + \left(b - \frac{\omega^2}{2\lambda}\right)t^2 \right].
		\end{align}
		To ensure $\mathcal{A} > 0$, we choose $t$ to minimize the polynomial in brackets. Setting $t^2 = a/2$, we get:
		\begin{align}
			\mathcal{A} = -\frac{\lambda (a/2)}{2} \left( \frac{a^2}{4} - \frac{a^2}{2} + b - \frac{\omega^2}{2\lambda} \right) = \frac{\lambda a}{4} \left( \frac{\omega^2}{2\lambda} - \left(b - \frac{a^2}{4}\right) \right).
		\end{align}
		By hypothesis \eqref{neccCond3}, $\mathcal{A} > 0$. The linear coefficient $\mathcal{B}$ collects the remaining positive terms (gradient and potential corrections):
		\begin{align}
			\mathcal{B} = \frac{t^2}{2} + \lambda \left[ t^6 - at^4 + \left(b - \frac{\omega^2}{2\lambda}\right)t^2 \right] + \lambda \left( \frac{a}{5}t^4 - \frac{1}{7}t^6 - \frac{1}{3}\left(b - \frac{\omega^2}{2\lambda}\right)t^2 \right),
		\end{align}
		and $\mathcal{C} = \frac{N^2 t^2}{2}$. For large $P$, the term $\mathcal{C} \ln P$ is sub-dominant to $\mathcal{B} P$. We can strictly bound $\ln P < P$ for $P>1$ to simplify the sufficient condition to a quadratic inequality:
		\begin{align}
			-\mathcal{A} P^2 + (\mathcal{B} + \mathcal{C}) P < 0 \implies P > \frac{\mathcal{B} + \mathcal{C}}{\mathcal{A}}.
		\end{align}
		Thus, for sufficiently large $P$, $I_{\omega}(\phi_0) < 0$.
		$\qquad\square$
	\end{proof}

\end{document}
